\documentclass[11pt]{article}
\usepackage[margin=0.95in]{geometry}
\usepackage{amsmath,amssymb}
\usepackage{graphicx}
\usepackage{booktabs,longtable,array}
\usepackage{xcolor}
\usepackage[hidelinks]{hyperref}
\usepackage{enumitem}
\usepackage{caption}
\setlist{itemsep=2pt,topsep=4pt}
\captionsetup{font=small,labelfont=bf}
\setlength{\parskip}{5pt}
\setlength{\parindent}{0pt}
\renewcommand{\arraystretch}{1.15}
\title{Stage D: Embedded Export Support}
\author{}
\date{trm-delay-robust-scaling $\cdot$ Phase 4 stage note $\cdot$ 2026-10-08\\[2pt]Plan \S{}9. Engineering only, no gate. Outputs in \texttt{.\allowbreak{}.\allowbreak{}/\allowbreak{}results/\allowbreak{}export\_\allowbreak{}c/\allowbreak{}}; script \texttt{.\allowbreak{}.\allowbreak{}/\allowbreak{}.\allowbreak{}.\allowbreak{}/\allowbreak{}experiments/\allowbreak{}phase4/\allowbreak{}export\_\allowbreak{}c.\allowbreak{}py}.}
\begin{document}
\maketitle
\section*{1. What was asked}



A script to export the actor weights of three Pendulum actors to float32 C headers for an on-device test on the author's STM32: MLP 135k, TRM 136k (H = 3) and TRM H = 1. For each, write 20 reference input/output vector pairs from PyTorch, and report the parameter count and the FP32 and INT8 weight sizes against a 512 KiB flash budget. The STM32 side is the author's.



\section*{2. What was done}



\texttt{export\_c.py} rebuilds each actor from its results file (Phase 3 Stage 1, fixed d = 1 + queue, seed 1) and writes one header per actor containing:



\begin{itemize}
\item every parameter tensor as a \texttt{static const float} array in PyTorch's row-major layout;
\item the trunk hyperparameters as \texttt{\#define}s, and for TRM the fixed recursion buffers (the initial latents \texttt{H\_init}, \texttt{L\_init}), which the forward pass needs but which are not trainable parameters;
\item 20 reference pairs: inputs are observations visited by the deterministic policy (4-dimensional: three plant observations and the one-slot action queue), outputs are the PyTorch float32 actions.
\end{itemize}



INT8 size assumes per-tensor symmetric quantisation: one byte per weight plus one float32 scale per tensor, with buffers kept in float32. No quantised model was built.



\begin{center}\small\begin{tabular}{>{\raggedright\arraybackslash}p{0.215\linewidth}>{\raggedleft\arraybackslash}p{0.115\linewidth}>{\raggedleft\arraybackslash}p{0.090\linewidth}>{\raggedleft\arraybackslash}p{0.166\linewidth}>{\raggedleft\arraybackslash}p{0.146\linewidth}>{\raggedright\arraybackslash}p{0.090\linewidth}}
\toprule
Actor (seed 1) & Actor parameters & Buffer floats & FP32 size & INT8 size & Fits 512 KiB \\
\midrule
MLP 135k (\texttt{dnn}) & 67,329 & 0 & 263.0 KiB (269,316 B) & 65.8 KiB (67,353 B) & FP32 and INT8 \\
TRM 136k (\texttt{trm\_small}, H = 3) & 67,805 & 88 & 265.2 KiB (271,572 B) & 66.6 KiB (68,185 B) & FP32 and INT8 \\
TRM H = 1 (\texttt{trm\_h1}, 788k config) & 393,857 & 256 & 1539.5 KiB (1,576,452 B) & 385.7 KiB (394,909 B) & INT8 only \\
\bottomrule
\end{tabular}\end{center}



\emph{Source: \texttt{phase4/\allowbreak{}results/\allowbreak{}export\_\allowbreak{}c/\allowbreak{}export\_\allowbreak{}summary.\allowbreak{}json}. Parameter counts here are the actor alone; the "135k"/"788k" labels are actor + critic totals under DDPG. TRM 788k (H = 3) has the same parameters as TRM H = 1 and would have the same size; it was not exported.}



\section*{3. Reading the table}



\begin{itemize}
\item \textbf{At matched size, the two trunks are the same flash cost.} MLP 135k and TRM 136k actors differ by 476 parameters and 88 buffer floats; both fit the budget in FP32 with room to spare.
\item \textbf{The 788k-configuration TRM fits only after quantisation.} In FP32 it needs about three times the budget; INT8 brings it to 75\% of it.
\item \textbf{Flash is not where TRM's embedded cost lies.} The TRM 136k actor still performs 21 cell applications per decision. Phase 3 measured it at 2.69 ms per decision (p99, one thread of a Katana Xeon) against the MLP's 0.064 ms, and TRM H = 1 at 1.42 ms. On a microcontroller both would be far slower; the ratio, not the size, is the deployment question.
\end{itemize}



\section*{4. Not done}



\begin{itemize}
\item \textbf{The on-device equivalence test} (the author's STM32 side) has not been run, so it is not known whether the headers reproduce PyTorch's outputs on the target to within float32 tolerance.
\item \textbf{Quantised accuracy.} The INT8 figures are sizes only. No INT8 actor was evaluated, so the effect of quantisation on return is unknown.
\item \textbf{On-device latency} was not measured.
\item Only seed 1 of each actor was exported.
\end{itemize}
\end{document}
